% !TeX program = xelatex
\documentclass[UTF8,fontset=fandol,10pt,a4paper]{ctexart}
\usepackage[left=22mm,right=22mm,top=23mm,bottom=22mm,headheight=14pt]{geometry}
\usepackage{fontspec}
\setmainfont{TeX Gyre Heros}
\setsansfont{TeX Gyre Heros}
\setmonofont{Latin Modern Mono}
\renewcommand{\familydefault}{\sfdefault}
\usepackage[table]{xcolor}
\usepackage{tikz}
\usetikzlibrary{arrows.meta,positioning,fit,backgrounds,calc}
\usepackage{booktabs,tabularx,array,enumitem,float,caption,fancyhdr,hyperref}
\definecolor{blue}{HTML}{2869A9}
\definecolor{ink}{HTML}{293340}
\definecolor{muted}{HTML}{647181}
\definecolor{line}{HTML}{C5D4E2}
\definecolor{pale}{HTML}{F0F5FA}
\hypersetup{colorlinks=true,urlcolor=blue,linkcolor=blue,pdftitle={Redflag 引擎架构},pdfauthor={Redflag}}
\color{ink}
\setlength{\parindent}{0pt}
\setlength{\parskip}{6pt}
\linespread{1.18}
\setlist[itemize]{leftmargin=1.5em,itemsep=2pt,topsep=3pt}
\renewcommand{\arraystretch}{1.45}
\newcolumntype{Y}{>{\raggedright\arraybackslash}X}
\newcolumntype{L}[1]{>{\raggedright\arraybackslash}p{#1}}
\ctexset{section={format=\Large\bfseries\color{blue},beforeskip=18pt,afterskip=8pt},subsection={format=\large\bfseries\color{blue},beforeskip=12pt,afterskip=5pt}}
\captionsetup{font=small,labelfont=bf,textfont={color=muted},skip=7pt}
\pagestyle{fancy}\fancyhf{}
\fancyhead[L]{\small\color{muted}Redflag · C++ / OpenGL · 技术架构}
\fancyhead[R]{\small\color{muted}2026-11-21}
\fancyfoot[C]{\small\color{muted}\thepage}
\renewcommand{\headrulewidth}{0pt}
\newcommand{\note}[1]{{\small\color{muted}#1\par}}
\newcommand{\code}[1]{{\small\texttt{#1}}}
\newcommand{\ruleline}{\par\vspace{-4pt}{\color{line}\hrule}\vspace{8pt}}
\tikzset{
 module/.style={draw=line,fill=pale,rounded corners=2pt,align=center,inner sep=8pt,text width=3.9cm,minimum height=1.6cm,font=\small},
 wide/.style={module,text width=10.2cm,minimum height=1.2cm},
 flow/.style={-{Stealth[length=2mm]},draw=muted,line width=.65pt},
 label/.style={font=\scriptsize,fill=white,inner sep=2pt,text=muted}}

\begin{document}
\zihao{-5}
{\Huge\bfseries Redflag 引擎架构}\par\vspace{5pt}
{\large\color{muted}从导入一个模型，到在编辑器里看到它}\par
{\color{blue}\hrule height 1pt}\vspace{10pt}
Redflag 是面向 Windows 的 C++ 游戏引擎，使用 OpenGL 绘制。RedflagForge 是编辑器：在这里摆放模型、调整灯光，或启动场景查看运行结果。本文沿着一个模型进入画面的过程，解释各部分为什么存在，以及代码从哪里开始读。

\section{模型从导入到显示}\ruleline
假设把一把椅子放进场景。引擎先记住它对应哪个文件，第一次需要绘制时再读出模型。之后每一帧使用已有模型，结合椅子的位置、相机和灯光生成画面。
\begin{figure}[H]\centering
\begin{tikzpicture}
\node[module] (a) at (0,0) {\textbf{1 记住模型身份}\\模型文件旁保存 ID\\场景记录这个 ID};
\node[module] (b) at (5.2,0) {\textbf{2 取出模型}\\已有就复用\\首次需要时读取文件};
\node[module] (c) at (10.4,0) {\textbf{3 提交这把椅子}\\模型的三角形\\位置、材质、实体编号};
\node[module] (d) at (0,-2.7) {\textbf{4 先画辅助图}\\记录遮挡与表面朝向\\供光照计算使用};
\node[module] (e) at (5.2,-2.7) {\textbf{5 计算主画面}\\结合材质和灯光\\画出可见表面};
\node[module] (f) at (10.4,-2.7) {\textbf{6 显示到编辑器}\\调整亮度与泛光\\把结果放进视口};
\draw[flow](a)--(b);\draw[flow](b)--(c);
\draw[flow](c.south)--++(0,-.45)-|(d.north);
\draw[flow](d)--(e);\draw[flow](e)--(f);
\end{tikzpicture}\caption{本文的阅读主线。辅助图用于计算，最终画面才交给编辑器显示。}\end{figure}

\subsection{资产与渲染图像的分工}
同一把椅子可以在场景里放十次，模型文件只需加载一份，每个物体另存自己的位置。磁盘上可重复使用的模型、贴图等文件，在本文中称为\textbf{资产}。它们的身份和加载由资产代码管理。

画面中的阴影、深度等结果会随相机或灯光变化，属于\textbf{渲染生成的图像}。它们由渲染代码创建和更新。这样修改视口大小时，只需重建相关图像，不必重新导入椅子。

\note{第 2、3 节解释谁发起这些操作；第 4 节跟踪模型加载；第 5 节跟踪一帧图像；第 6、7 节用于查代码和理解释放顺序。第 8 节单列尚未实现的工作。}

\clearpage
\section{模块职责}\ruleline
编辑器保存“场景里有哪些物体”，场景渲染代码把这些物体转成绘制工作，OpenGL 实现再把工作交给显卡。图中先写用途，下面保留对应类名，方便沿着源码查找。
\begin{figure}[H]\centering
\begin{tikzpicture}
\node[module] (forge) at (0,0) {\textbf{编辑场景与显示结果}\\RedflagForge\\ForgeLayer / Panels};
\node[module] (scene) at (5.2,0) {\textbf{保存和更新物体}\\Scene\\位置、模型、灯光、脚本};
\node[module] (sr) at (10.4,0) {\textbf{安排渲染顺序}\\SceneRenderer\\收集物体与灯光};
\node[module] (app) at (0,-3.0) {\textbf{维持窗口与主循环}\\Application\\Window / LayerStack};
\node[module] (assets) at (5.2,-3.0) {\textbf{找到并复用模型}\\AssetLibrary\\AssetManager};
\node[module] (draw) at (10.4,-3.0) {\textbf{组织几何绘制}\\Renderer2D / Renderer3D\\合并或逐次提交};
\node[wide] (gl) at (5.2,-5.8) {\textbf{在显卡上创建资源并执行绘制}\\Platform / OpenGL：缓冲、贴图、绘制程序和绘制目标的实现};
\draw[flow](app)--(forge);\draw[flow](forge)--(scene);\draw[flow](scene)--(sr);\draw[flow](sr)--(draw);
\draw[flow,dashed](scene)--node[label,right]{查找模型}(assets);
\draw[flow](app.south)|-(gl.west);\draw[flow](draw.south)|-(gl.east);
\end{tikzpicture}\caption{实线表示主要调用方向，虚线表示资源查询。}\end{figure}
\subsection{以移动椅子为例}
在编辑器里拖动椅子，改变的是场景中这个物体的位置。下一帧，场景渲染代码取出新位置，把它与模型一起提交。模型的三角形和贴图仍然可以复用，不需要因为位置变化就重新读文件。

这里的“物体”在代码中叫 \code{Entity}。位置、模型和脚本等信息分别放在组件中，由 \code{Scene} 保存。EnTT 是负责存储这些实体与组件的外部库。

\note{图表示目前的主要分工，并不表示代码已经完全分开。Scene 的更新函数仍会触发绘制，一些编辑面板还放在 Scene/Panel 中，高层渲染函数也仍有直接 OpenGL 调用。}

\clearpage
\section{帧更新与场景状态}\ruleline
窗口打开后，应用不断更新和绘制。相邻两次更新之间的时间用于推进物体运动。编辑器根据当前模式，选择要更新哪个场景、用哪台相机看它。
\begin{figure}[H]\centering
\begin{tikzpicture}[node distance=8mm]
\node[wide] (a) {\textbf{1 开始这一帧}\\Application::Run 计算经过的时间，更新各个 Layer};
\node[wide,below=of a] (b) {\textbf{2 选择场景和相机}\\ForgeLayer::OnUpdate 准备接收画面的目标};
\node[wide,below=of b] (c) {\textbf{3 更新物体}\\Scene::OnUpdate* 按模式执行脚本或物理，再进入绘制};
\node[wide,below=of c] (d) {\textbf{4 生成场景画面}\\SceneRenderer 收集物体与灯光，按顺序完成计算};
\node[wide,below=of d] (e) {\textbf{5 显示画面与编辑工具}\\ForgeLayer 把结果交给 ImGui，显示视口和面板};
\draw[flow](a)--(b);\draw[flow](b)--(c);\draw[flow](c)--(d);\draw[flow](d)--(e);
\end{tikzpicture}\caption{Layer 是应用更新的功能层，编辑器的主要工作放在 ForgeLayer 中。}\end{figure}
\subsection{编辑场景与运行副本}
\par
\begin{tabularx}{\linewidth}{@{}L{2.2cm}Y Y@{}}\toprule
模式 & 场景和相机 & 这一帧做什么\\\midrule
编辑 Edit & 原编辑场景；编辑器相机 & 编辑物体并绘制。\\
播放 Play & 编辑场景的副本；场景主相机 & 运行脚本、物理和粒子，再绘制。\\
模拟 Simulate & 编辑场景的副本；编辑器相机 & 推进物理并绘制。\\\bottomrule
\end{tabularx}\par\vspace{4pt}
播放和模拟使用副本，因此运行中的位置变化不会直接写回编辑场景。停止时，先结束脚本或物理，再丢弃副本并恢复编辑场景。悬停对象等引用也要清除，否则还可能指向已经销毁的运行副本。

\clearpage\section{模型编号与文件位置}\ruleline
把椅子模型拖入属性面板时，编辑器先调用 \code{Add<Model>} 登记这份文件，再调用 \code{Load<Model>}。模型读成功后，组件才记下路径、编号和模型对象，视口可以立刻显示结果。下面的图把这两个动作分开：登记时只记录模型是谁、在哪里，加载时才会读出三角形、材质和贴图。

\begin{figure}[H]\centering
\begin{tikzpicture}
\node[module] (source) at (0,0) {\textbf{1 文件旁保存编号}\\模型 + .rfmeta\\编号随元数据保留};
\node[module] (library) at (5.2,0) {\textbf{2 记录编号对应的路径}\\AssetLibrary\\只登记和查询};
\node[module] (manager) at (10.4,0) {\textbf{3 请求模型}\\AssetManager::Load\\先检查是否已加载};
\node[module] (component) at (0,-3) {\textbf{场景中的模型组件}\\保存 ModelID\\也保留旧路径迁移信息};
\node[module] (load) at (5.2,-3) {\textbf{4 首次加载时读文件}\\Model::Load\\解析几何、材质和贴图};
\node[module] (gpu) at (10.4,-3) {\textbf{5 返回可用的模型}\\创建显卡资源\\缓存与场景共同持有};
\draw[flow](source)--(library);\draw[flow](library)--(manager);
\draw[flow](component.north east)--node[label,left]{按编号找位置}(library.south west);
\draw[flow](manager.south)--node[label,right]{已有就复用}(gpu.north);
\draw[flow](manager.south west)--node[label,above left]{尚未加载}(load.north east);
\draw[flow](load)--(gpu);
\end{tikzpicture}\caption{登记只记录文件信息。读取模型和创建显卡资源发生在加载时。}\end{figure}
\subsection{模型改名与场景重开}
模型旁的 \code{.rfmeta} 文件保留一个 64 位无符号编号，代码中叫 \code{AssetID}，0 表示没有指定资源。场景保存这个编号。假如把 \code{chair.obj} 移到 \code{assets} 里的另一个文件夹，连同它的元数据一起移动，模型就能沿用原编号。编辑器的文件面板隐藏元数据，因此拖入场景时看见的是模型本身。

重新打开场景走的是另一条路径。\code{SceneSerializer::Deserialize} 先让 \code{AssetLibrary::Scan} 扫描 \code{assets}，从找到的元数据重建编号与路径的对应关系。遇到只有路径的旧场景，读取组件时按原路径登记并补上编号。对于扫描范围外的文件，保存的路径仍须可用，且元数据中的编号要与场景一致。文件移出这个范围后，仅凭原来的编号无法遍历磁盘找回它。

\subsection{网格和材质的复用}
场景重开后，组件记得模型编号，可以等到绘制第一次用到它时再请求加载。\code{AssetLibrary::Get} 回答文件的类型和位置，\code{AssetManager::Load} 按类型和编号查询缓存。缓存里如已有这把椅子，就直接返回同一份模型；首次加载时才通过 Assimp 读取文件并创建子网格和材质。一个子网格能独立绘制，例如椅背与椅腿各自用一种材质。

目前，解析文件和创建显卡资源都在调用加载的线程上完成，因此首次读取大模型可能让这一帧等一会儿。后台加载、内存预算和文件改动后的自动重载均属于后续设计。模型用完后，\code{Unload} 或 \code{ClearAll} 可以撤掉缓存的引用；场景仍持有的模型会继续使用。

\clearpage\section{场景颜色与后处理}\ruleline
看视口里的椅子，椅背有自己的颜色，椅腿会投影到地面，调高曝光后整幅画面又一起变亮。这几件事发生在不同的绘制阶段。开始时，\code{DrawModel} 把各子网格、位置、材质和实体编号记下来。\code{SceneRenderer} 按下图的次序使用它们。主画面把同材质的记录排在一起；共享网格和同组贴图的不透明物体也可用一次实例化绘制提交，仍能分别选中。

\begin{figure}[H]\centering
\begin{tikzpicture}[module/.append style={text width=5.6cm,minimum height=1.35cm,inner sep=6pt}]
\node[module] (a) at (0,0) {\textbf{1 画二维物体，收集三维物体}\\同时取出这一帧的灯光数据};
\node[module] (b) at (7.8,0) {\textbf{2 从灯光位置看场景}\\Shadow 记遮挡；RSM 记表面反射信息};
\node[module] (c) at (0,-2.25) {\textbf{3 从相机位置记录表面}\\深度记远近；法线记表面朝向};
\node[module] (d) at (7.8,-2.25) {\textbf{4 用辅助图估算光照}\\SSAO 算局部遮蔽；RSM 算间接光};
\node[module] (e) at (0,-4.5) {\textbf{5 混合本帧与前帧间接光}\\RSM History 减少采样结果的闪动};
\node[module] (f) at (7.8,-4.5) {\textbf{6 绘制模型的表面颜色}\\Renderer3D::EndScene 使用材质和光照};
\node[module] (g) at (0,-6.75) {\textbf{7 补上天空盒与调试线框}\\结果仍保存到场景的颜色图中};
\node[module] (h) at (7.8,-6.75) {\textbf{8 调整亮度，加入泛光}\\PostProcess 输出供编辑器显示的图像};
\draw[flow](a)--(b);\draw[flow](c)--(d);\draw[flow](e)--(f);\draw[flow](g)--(h);
\draw[flow](b.south)--++(0,-.35)-|(c.north);
\draw[flow](d.south)--++(0,-.35)-|(e.north);
\draw[flow](f.south)--++(0,-.35)-|(g.north);
\end{tikzpicture}\caption{当前主要顺序。是否执行某项效果，还取决于开关、灯光与所需图像是否可用。}\end{figure}
\subsection{阴影、深度与表面颜色}
先从方向光望向场景，记下光线遇到的最近表面，地面着色时便能判断椅腿是否挡住灯光。打开反射阴影图（RSM）时，还会记录灯所照见的表面颜色等信息，供后面估算墙面反射的光。换回主相机后，深度给出每个屏幕位置可见的表面，法线给出这个表面的朝向。屏幕空间环境遮蔽（SSAO）比较邻近表面，估计椅腿与地面接触处能收到多少环境光。

间接光计算把眼前的表面位置投回上一帧，用先前的颜色减少闪动。投到画面外、相机后方或裁剪范围外的点会跳过；切换场景或修改投影、光照及采样设置也会清掉旧图。运动物体和刚露出的表面仍可能留下拖影，因为当前没有保存物体的运动和上一帧对应的深度。同投影的运行时相机互换也尚不能单独识别。

\subsection{曝光、泛光与编辑器视口}
\code{ForgeLayer} 准备一个 \code{Framebuffer} 接收场景颜色、实体编号、观察空间法线和深度。颜色先写成浮点数，灯或天空的亮度可以超过 1。开启泛光时，\code{PostProcess::Render} 从这张高动态范围（HDR）图提取亮部，再把不同尺寸的亮部图像合成柔和的光晕。泛光默认用六层，可以设置为一到六层。最后，后处理按照曝光把亮度映射到显示范围，编辑器将最终图像交给 ImGui 的视口面板。

目前 \code{SceneRenderer} 持有一个 \code{PostProcess} 来管理曝光、泛光和输出图。其他效果仍有静态共享数据，两个视口的图像尚未全部分开。调试面板里能看到的深度预览，是同步读回整张深度图后拉伸得到的灰度。若要从灰度求真实距离，还要用写入深度时的相机投影还原。


\clearpage
\section{源码入口与数据类型}\ruleline
按操作找入口，比先记住全部类名更容易。下面保留函数名，说明进入后应当关注什么。
\par
\begin{tabularx}{\linewidth}{@{}L{3.0cm}L{4.2cm}Y@{}}\toprule
功能 & 入口 & 查看内容\\\midrule
持续更新 & Application::Run & Layer 的更新、事件传递和退出。\\
场景读写 & SceneSerializer & Serialize / Deserialize 如何读写组件和模型 ID。\\
播放与停止 & Scene 的运行入口 & 复制场景后怎样启动脚本、物理，以及何时停止。\\
模型加载 & AssetManager::Load & 先查共享缓存，缺少时调用 Model::Load。\\
绘制记录 & Renderer3D::DrawModel & 逐个子网格组合变换，记录材质和实体编号。\\
阶段顺序 & SceneRenderer & RenderScene 用于编辑视图，RenderRuntime 用于运行相机。\\
后处理 & PostProcess::Render & 读取主画面，生成泛光并调整显示亮度。\\\bottomrule
\end{tabularx}\par\vspace{4pt}
\subsection{网格、材质与绘制记录}
模型里的三角形由 \code{Mesh} 保存。表面颜色、粗糙度和贴图引用由 \code{Material} 保存。\code{DrawCommand} 是本帧要画一次模型的记录，把网格、位置变换、材质和实体编号放在一起。材质可以复用，同一个网格也可以出现在多条绘制记录中。

运行在显卡上、负责计算顶点位置和像素颜色的小程序称为\textbf{着色器}，代码对象叫 \code{Shader}。贴图、网格数据和着色器配合起来，才足以完成一次绘制。

\subsection{外部库职责}
\par
\begin{tabularx}{\linewidth}{@{}L{4.6cm}Y@{}}\toprule
库 & 用途\\\midrule
GLFW / ImGui / ImGuizmo & 建立窗口和输入；绘制编辑面板；提供移动、旋转等操作工具。\\
Assimp / stb\_image / tinyexr & 读取模型文件；解码普通图像和 EXR 图像。\\
Mono / yaml-cpp & 执行 C\# 脚本；读写场景文本。ScriptGlue 连接脚本与 C++。\\
EnTT / Box2D / GLM / spdlog & 存储实体组件；计算二维物理；提供向量矩阵；记录日志。\\\bottomrule
\end{tabularx}\par\vspace{4pt}
\note{当前高层代码仍直接调用部分外部库，是否继续封装需要结合具体使用范围决定。}

\clearpage
\section{生命周期与资源释放}\ruleline
显卡上的贴图和缓冲需要通过 OpenGL 删除。执行删除时，窗口对应的图形上下文必须仍有效。这里的上下文就是执行 OpenGL 操作所需的环境。因此退出顺序从使用资源的场景开始，最后才关闭窗口。
\begin{figure}[H]\centering
\begin{tikzpicture}[node distance=7mm]
\node[wide] (a) {\textbf{1 停止场景和编辑器}\\LayerStack::Clear；ForgeLayer 退出播放或模拟，释放场景引用};
\node[wide,below=of a] (b) {\textbf{2 关闭脚本环境}\\ScriptEngine::Shutdown；释放托管运行环境和脚本数据};
\node[wide,below=of b] (c) {\textbf{3 清空资产缓存}\\AssetManager::ClearAll；释放缓存引用，再清空资源库};
\node[wide,below=of c] (d) {\textbf{4 关闭渲染器}\\Renderer::Shutdown；释放后处理、三维和二维绘制资源};
\node[wide,below=of d] (e) {\textbf{5 最后关闭窗口}\\释放 GraphicsContext，销毁 GLFW 窗口};
\draw[flow](a)--(b);\draw[flow](b)--(c);\draw[flow](c)--(d);\draw[flow](d)--(e);
\end{tikzpicture}\caption{先释放使用者和资源，再结束 OpenGL 的执行环境。}\end{figure}
\subsection{共享引用与卸载}
场景和缓存可以同时持有同一个模型。共享引用在代码中写作 \code{Ref<T>}，它是 \code{std::shared\_ptr} 的别名。\code{Unload} 只放开缓存的那份引用。场景还在使用时，模型继续存在。最后一份引用释放后，模型及其持有的资源才有机会销毁。

编辑器的“卸载模型”还会清空组件里的模型 ID 和路径，避免下一帧又按旧记录加载回来。阴影图和上一帧图像由渲染模块持有，不放进模型缓存。

\note{Layer 在栈内时由 LayerStack 管理。Clear 会结束并销毁全部 Layer；Pop 只调用 OnDetach 后移出，调用方接手对象，负责后续释放。}

\clearpage
\section{后续改造与验收}\ruleline
前七节说明当前代码。下表是尚未完成的改造，每一步都围绕一个实际问题，并给出完成后可以检查的结果。
\par
\begin{tabularx}{\linewidth}{@{}L{3.5cm}Y Y@{}}\toprule
现在的限制 & 下一步修改 & 验收结果\\\midrule
读模型与上传显卡混在一起 & 先把文件解析和显卡资源创建分开，再明确导入设置与保存产物。 & 修改设置后重新导入，运行时能使用对应产物。\\
场景更新里直接触发绘制 & 先收集本帧需要的数据，再交给渲染代码；拟用 RenderSceneData 保存。 & 更新和绘制能分别调用，绘制不改场景组件。\\
多个视口会共用部分图像 & 让每个视口持有自己的尺寸、效果设置和前帧图像。 & 两台相机使用不同尺寸连续绘制，结果不相互影响。\\
效果依赖写在大函数中 & 把各步骤的输入图、输出图和绘制状态写清楚，仍按固定顺序执行。 & 开关效果、调整大小后，结果与原流程一致。\\
临时图像还由各处管理 & 前面步骤稳定后，再评估是否让 RenderGraph 根据图像依赖安排执行和复用。 & 有实测的资源收益，且画面保持一致。\\\bottomrule
\end{tabularx}\par\vspace{4pt}
\subsection{目前已经验证的范围}
资产测试检查稳定 ID、文件移动和共享引用。后处理测试在隐藏 OpenGL 4.6 窗口中绘制，验证两个实例使用不同曝光、交替绘制、调整尺寸和切换泛光。Shader 缓存、模型导入及透明裁剪轮廓另有专项测试。具体输入与结果见《Redflag 调试记录》。这些测试不能代替完整编辑器交互回归。

图形接口与 Vulkan 的迁移单列于 RHI 文档。RHI 是隔开渲染代码与具体图形 API 的接口，本项目当前仍使用 OpenGL。

\subsection{参考来源与采用范围}
\href{https://github.com/TheCherno/Hazel/blob/master/Hazel/src/Hazel/Renderer/Shader.h}{Hazel 的 ShaderLibrary} 使用 Add / Load / Get / Exists 管理 Shader。Redflag 沿用这种易懂的命名。本项目的 AssetID 和 AssetLibrary 是自己的扩展，不能据此推断 Hazel 公开版有相同的资产设计。

\href{https://docs.unity.com/en-us/engine/6000.3/manual/assets-and-media/asset-database/contents}{Unity 6.3 的 Asset Database} 将源文件、元数据和导入结果联系起来，帮助理解为什么身份要与路径分开。\href{https://github.com/Unity-Technologies/Graphics/blob/master/Packages/com.unity.render-pipelines.core/Documentation~/render-graph-system.md}{Unity 的 Render Graph 文档} 用于后续研究图像依赖与资源复用，当前没有照搬完整系统。
\end{document}
